\section{Scheduling the Analysis}
\label{sec:method}
\subsection{Frames and processing calls}
Let $x[n]$ be a real signal sampled at $f_s$\,Hz. The analyzer uses a
power-of-two transform length $N$ and an integer hop $H\ge1$, the number of
samples between successive frame endpoints. Frame $r$ ends at $t_r=rH$ and
contains
\begin{equation}
 u_r[m]=x[t_r-N+1+m]w[m],\quad 0\le m<N,
 \label{eq:frame}
\end{equation}
where $w[m]$ is the selected window, including any fixed gain normalization.
Its spectrum is
\begin{equation}
 X_r[k]=\sum_{m=0}^{N-1}u_r[m]e^{-j2\pi km/N}.
 \label{eq:dft}
\end{equation}
Deferring computation changes when this result is available, not the samples
in the sum. In particular, later input must not overwrite a sample before the
scheduled windowing operation has read it.

A \emph{sample call} advances the engine by one sample. A \emph{block}
contains $D$ such calls and is the unit delivered by the audio host. The hop
and block size need not match or divide one another. At 48\,kHz, a block of
64 samples has a time budget of 1.333\,ms; a hop of 1024 samples gives
21.33\,ms between spectra. These are different clocks: distributing a frame
over its hop can move work out of one block and into subsequent blocks.

\begin{figure}[tb]
\centering
\begin{tikzpicture}[x=0.57cm,y=0.6cm,font=\footnotesize,>=Latex]
\node[anchor=west] at (0,3.1) {\textbf{Batch: one long processing call}};
\draw[->] (0,1.85) -- (13.8,1.85);
\draw[fill=accent!25,draw=accent] (0.2,1.85) rectangle (1.1,2.85);
\node[anchor=west] at (1.5,2.45) {window $+$ FFT $+$ output};
\fill[accent] (0.65,1.85) circle (2pt);
\node[anchor=west] at (1.5,2.02) {publish immediately};
\node[anchor=west] at (0,0.95) {\textbf{Distributed: resume on successive calls}};
\draw[->] (0,-0.2) -- (13.8,-0.2);
\foreach \x in {0.2,1.2,2.2} {
  \draw[fill=ink!20,draw=ink] (\x,-0.2) rectangle +(0.65,0.48);
}
\foreach \x in {3.2,4.2,5.2,6.2,7.2,8.2,9.2} {
  \draw[fill=ink!65,draw=ink] (\x,-0.2) rectangle +(0.65,0.48);
}
\foreach \x in {10.2,11.2,12.2} {
  \draw[fill=accent!45,draw=accent] (\x,-0.2) rectangle +(0.65,0.48);
}
\node[anchor=north,align=center] at (1.5,-0.3) {window /\\pack};
\node[anchor=north] at (6.5,-0.3) {FFT butterflies};
\node[anchor=north,align=center] at (11.5,-0.3) {magnitudes /\\output};
\draw[<->] (0.2,-1.5) -- (12.85,-1.5);
\node[fill=white,inner sep=2pt] at (6.5,-1.5) {one hop: $H$ sample calls};
\draw[->,accent] (12.85,0.85) -- (12.85,0.3);
\node[anchor=east] at (12.65,0.7) {publish};
\end{tikzpicture}
\caption{The same selected frame can be computed at its endpoint or spread
over the following hop. Input capture continues in both cases. Bars are
conceptual work slices, not timing measurements; stages and bar heights are
not to scale. The distributed result is published on the last call.}
\label{fig:schedule}
\end{figure}


\subsection{A resumable FFT}
An iterative radix-2 FFT traverses stages, groups within a stage, and pairs
within a group. A butterfly reads two coefficients, multiplies one by a
precomputed twiddle factor, and replaces the pair by their sum and difference.
Saving the current stage, group, and pair is enough to resume the traversal.
The coefficient array remains in place between calls. Pausing introduces no
new transform mathematics: when resumed, the next butterfly sees the same
operands that an uninterrupted traversal would have produced.

For real input, pair the windowed samples as
$z[m]=u_r[2m]+j u_r[2m+1]$ and compute an $M=N/2$ point complex FFT.
Standard real-spectrum reconstruction then produces the $K=M+1$ bins from
DC through Nyquist~\cite{sorensen1987}. The complex traversal contains
\begin{equation}
 B=\frac{M}{2}\log_2 M=\frac{N}{4}\log_2(N/2)
 \label{eq:butterflies}
\end{equation}
butterflies. The traversal indices require constant storage; the coefficient
and lookup arrays require $O(N)$ storage. Any partition of this ordered
traversal preserves the result within the same arithmetic implementation.
This does not imply bitwise agreement with a different FFT library.

\subsection{Include the work around the FFT}
The analyzer divides each frame into four stages
(Table~\ref{tab:stages}). First, it windows pairs of retained input samples
and stores them in the FFT's permuted order. Next, it executes the
butterflies. It then reconstructs positive-frequency bins and computes their
magnitudes. Finally, it applies enabled smoothing and emits display data one
bin at a time. Each stage finishes before the next begins; a sample call can
cross a stage boundary if it has work left to do.

\begin{table}[tb]
\centering\small
\begin{tabular}{@{}p{0.22\linewidth}p{0.55\linewidth}r@{}}
\toprule
Stage & One operation & Credits \\
\midrule
Prepare & Window, pack, and permute an input pair & $pM$ \\
Transform & Execute one complex butterfly & $B$ \\
Magnitude & Reconstruct one bin; compute magnitude and optional prefix sum & $K$ \\
Output & Smooth one bin and write its display data & $oK$ \\
\bottomrule
\end{tabular}
\caption{Scheduled stages for real input. $M=N/2$ pairs and $K=M+1$ bins
are processed in order. The last column gives the total credits per stage;
$p$ and $o$ are preparation and output weights.}
\label{tab:stages}
\end{table}


Frequency smoothing uses prefix sums of the magnitudes. Once those sums are
complete, the mean over any contiguous bin range needs two prefix reads and
a subtraction rather than a loop over that range. Temporal smoothing updates
one stored value per bin. Thus neither form of smoothing adds an unscheduled
whole-spectrum pass on the publication call. The display callback must also
perform bounded work per bin; an arbitrary expensive callback would defeat
the schedule.

To account for operations with different costs, the implementation assigns
integer \emph{credits}. Most operations use one credit. Rebuilding window
coefficients uses $p=4$ credits per input pair, compared with $p=1$ when the
cache is ready. Output uses $o=2$ credits per bin for Fourier's four-channel
display conversion and $o=1$ otherwise. A weighted operation executes at its
first credit; the remaining credits defer later work. These are engineering
weights, not measurements of equal execution time. The total is
\begin{equation}
 W=pM+B+(1+o)K.
 \label{eq:work}
\end{equation}
The weights affect where work falls within the hop without changing the
computed spectrum or the chosen publication time.

\subsection{An exact completion schedule}
Number a frame's sample calls $s=1,\ldots,H$, starting with the call that
captures its endpoint. Assign call $s$ the quota
\begin{equation}
 q_s=\floor{\frac{sW}{H}}-\floor{\frac{(s-1)W}{H}}.
 \label{eq:quota}
\end{equation}
The cumulative count after $s$ calls is $\floor{sW/H}$, so all $W$ credits
have been consumed after exactly $H$ calls. Every call receives either
$\floor{W/H}$ or $\ceil{W/H}$ credits. An integer quotient and remainder
implement the rule without dividing on every sample:
\begin{lstlisting}
// Once per frame configuration:
base = W / H; remainder = W % H;
// At frame start: error = 0.
// On each sample call:
quota = base;
error += remainder;
if (error >= H) { error -= H; ++quota; }
advance_analysis(quota);
if (++phase == H) { publish(); phase = 0; }
\end{lstlisting}
Here \code{advance\_analysis} consumes credits in the stage order above,
preserving the current stage and its cursor. Input capture occurs on every
call, including calls whose quota is zero. Publication waits for the fixed
boundary even if the last weighted operation has already executed.

For a clean scalar frame with $N=4096$ and $H=1024$, $W=17410$:
1022 calls consume 17 credits and two consume 18. For fixed configuration,
any 64 consecutive calls consume at most
$\ceil{64W/H}=1089$ credits, including blocks that cross a frame boundary.
More generally, when $H$ is proportional to $N$, the maximum per-sample
quota grows as $O(\log N)$ while a batch frame still requires
$O(N\log N)$ work at once. The schedule bounds operation counts, not
elapsed time: instruction costs, memory effects, and operating-system
interruptions are outside that arithmetic statement.

Keeping the frame clock separate from arithmetic completion is essential.
A fixed rounded quota such as $\ceil{B/H}$ butterflies per call can finish
early. Restarting immediately would change the actual update rate and any
smoothing calibrated to it. Equation~\eqref{eq:quota} distributes the
remainder across the intended interval and retains the requested cadence.

\subsection{Input lifetime and result delay}
The frame selected at $t_r$ is published at $t_r+H-1$. Measured from its
newest sample, the delay is therefore
\begin{equation}
 A_{\mathrm{endpoint}}=(H-1)/f_s.
 \label{eq:age}
\end{equation}
Measured from the window center, add $(N-1)/(2f_s)$. For the 4096-point,
1024-hop example at 48\,kHz, these ages are 21.31 and 63.97\,ms,
respectively. Even an immediate batch transform has the window-center age;
distribution adds the time spent advancing the frame. These are logical
sample-clock ages, excluding computation time, audio buffering, and repaint.

A preallocated input ring with capacity $N_{\max}+H_{\max}$ retains every
selected sample until analysis finishes. At most $H-1$ later samples arrive
before publication, so the oldest sample of the $N$-sample frame cannot be
overwritten. Windowing can therefore proceed incrementally without a bulk
frame copy. The first frames use zero padding for input that predates startup.
With fixed settings, only one frame is in progress and successive results are
published exactly $H$ samples apart.
